\documentclass[11pt]{amsart}
\usepackage[margin=1in]{geometry}
\usepackage{amssymb,amsmath,amsthm,mathtools}
\usepackage{booktabs,array,tabularx}
\usepackage{xcolor}
\usepackage[hidelinks]{hyperref}
\hypersetup{pdftitle={Machine-Checked Flow-Certificate Thresholds for Planar Graphs Using Rotation-System Certificates},pdfauthor={Filip Majchrzak}}
\usepackage{listings}
\IfFileExists{microtype.sty}{\usepackage{microtype}}{}
\emergencystretch=3em

\lstdefinelanguage{lean}{
  keywords={theorem,lemma,def,structure,where,by,fun,exact,refine,rw,obtain,namespace,end,open,import,variable,abbrev,noncomputable},
  sensitive=true, comment=[l]{--}, morecomment=[s]{/-}{-/}, string=[b]",
}
\lstset{language=lean, basicstyle=\ttfamily\footnotesize, columns=fullflexible, keepspaces=true,
  breaklines=true, frame=single, framerule=0.3pt, xleftmargin=2pt, xrightmargin=2pt,
  literate={ℕ}{{$\mathbb{N}$}}1 {ℝ}{{$\mathbb{R}$}}1 {ℤ}{{$\mathbb{Z}$}}1 {≠}{{$\neq$}}1 {∧}{{$\wedge$}}1
           {→}{{$\to$}}1 {∀}{{$\forall$}}1 {∃}{{$\exists$}}1 {ε}{{$\varepsilon$}}1 {≤}{{$\le$}}1
           {⟨}{{$\langle$}}1 {⟩}{{$\rangle$}}1 {↔}{{$\leftrightarrow$}}1 {ᵀ}{{$^{\mathsf T}$}}1 {ᵥ}{{$_v$}}1
           {∘}{{$\circ$}}1 {−}{{$-$}}1 {φ}{{$\varphi$}}1 {α}{{$\alpha$}}1 {·}{{$\cdot$}}1 {×}{{$\times$}}1
           {^[}{{\^{}[}}2 {|}{{$|$}}1
}

\theoremstyle{plain}
\newtheorem{theorem}{Theorem}
\newtheorem{proposition}[theorem]{Proposition}
\newtheorem{lemma}[theorem]{Lemma}
\newtheorem{corollary}[theorem]{Corollary}
\theoremstyle{definition}
\newtheorem{definition}{Definition}
\newtheorem{example}{Example}
\theoremstyle{remark}
\newtheorem{remark}{Remark}

\newcommand{\R}{\mathbb{R}}
\newcommand{\Z}{\mathbb{Z}}
\newcommand{\T}{\mathsf{T}}
\newcommand{\src}{\operatorname{src}}
\newcommand{\tgt}{\operatorname{tgt}}
\newcommand{\opt}{\operatorname{opt}}
\newcommand{\cost}{\operatorname{cost}}
\newcommand{\capac}{\operatorname{cap}}
\newcommand{\row}{\operatorname{row}}
\newcommand{\face}{\operatorname{face}}
\newcommand{\rot}{\operatorname{rot}}
\newcommand{\Adm}{\operatorname{Adm}}
\newcommand{\code}[1]{\texttt{#1}}

\title[Machine-checked flow-certificate thresholds for planar graphs]{Machine-Checked Flow-Certificate Thresholds\\ for Planar Graphs Using Rotation-System Certificates}
\author{Filip Majchrzak}
\address{Independent researcher, Poland (SAPIO MORPHEUS project)}
\thanks{ORCID 0009-0007-3475-8900. Version 1.1, 24 November 2026. Version 1.1 adds the literature on squared rectangles (Section~7 and Appendix~A); the results and the Lean code are unchanged. Archived at Zenodo: this version \href{https://doi.org/10.5281/zenodo.24496271}{doi:10.5281/zenodo.24496271}; all versions \href{https://doi.org/10.5281/zenodo.24488175}{doi:10.5281/zenodo.24488175}. Text licensed CC BY 4.0; Lean code Apache-2.0.}
\date{24 November 2026}

\begin{document}

\begin{abstract}
Let $G$ be a connected planar multigraph with unit resistances, let $s,t$ be vertices and let $e$ be an edge that
carries a nonzero current $I_e$ when one unit of current is sent from $s$ to $t$. From $G$ and its plane dual we
build two capacitated minimum-cost circulation networks on one shared topology, whose observed arcs are indexed by
the edges of $G$. We prove that the least common additive tolerance $\varepsilon$ for which both networks have
feasible circulations that agree on every observed arc, each within $\varepsilon$ of its own optimum, is
\[
\Delta^* \;=\; \frac{|I_e|}{1+|I_e|} \;=\; \frac{|N|}{D+|N|},
\]
where $D=\det(AA^{\T})>0$ is the reduced Laplacian determinant and $N\in\Z$ is a Cramer numerator with $I_e=N/D$.
Planarity enters only through the identity $\ker B=\row A$ between the reduced incidence matrices of $G$ and its
dual. The statement and its proof are formalised in Lean~4 with Mathlib. Planarity is represented by a genus-$0$
rotation-system certificate, from which the dual graph, the orthogonality $BA^{\T}=0$, the connectivity of the dual
and the identity $\ker B=\row A$ are derived inside Lean. The development has no \code{sorry} and no added axioms,
and every module is replayed by the kernel checker. We state precisely which ingredients are classical, which are
cited and not formalised (the equivalence between planarity and genus-$0$ rotation systems, and the matrix-tree
identification of $D$ and $|N|$), and which may be new. We make no claim of priority.
\end{abstract}

\maketitle

\section{Introduction}

Two optimisation problems often share observable quantities while differing in costs and capacities. One may then
ask how much optimality each must give up so that both report the \emph{same} observation. This paper answers that
question exactly for one specific pair of circulation networks attached to a planar electrical network. We call
a feasible circulation of either network a \emph{flow certificate}. The \emph{threshold} of the pair is the least
common optimality tolerance at which flow certificates of both networks can show the same observation.

The answer is governed by a single number from classical circuit theory, the current $I_e$ through an edge $e$
under a unit $s$--$t$ feed. The least shared tolerance is $|I_e|/(1+|I_e|)$. By Kirchhoff's theorem, $|I_e|$ is a
ratio of a signed spanning-tree count to the number of spanning trees. The threshold therefore has the integer
form $\kappa/(\tau(G)+\kappa)$, where $\tau(G)$ is the number of spanning trees and $\kappa$ is a signed count of
spanning trees.

The mathematics is elementary. It is a linear-algebraic argument that combines Kirchhoff's laws, planar duality
and the minimum of two linear functions. The main contribution is a machine-checked proof in Lean~4 of the certified form of
the theorem, for graphs given together with a genus-$0$ rotation-system certificate. The proof starts from this
concrete combinatorial representation of planarity and reaches the exact threshold on explicitly constructed
networks. The link between certificates and geometric planarity is classical and is cited, not formalised. At the
pinned Mathlib commit we found no theory of planar graphs, rotation systems or the matrix-tree theorem (by
searching the source), so the development builds the required pieces directly.

\subsection*{Contributions}
\begin{enumerate}
\item \textbf{The exact threshold.} The precise statement (Theorem~\ref{thm:main}) and a self-contained proof
  (Section~\ref{sec:proof}). The proof uses only connectivity, the bound $0\le I\le 1$ after reorientation, and
  the identity $\ker B=\row A$.
\item \textbf{Rotation-system certificates.} We derive $\ker B=\row A$ from a genus-$0$ rotation-system
  certificate (Section~\ref{sec:cert}), with the dual computed from the certificate.
\item \textbf{A Lean~4 formalisation.} About $1.7$ thousand lines, kernel-checked, with exact versions and a
  reproduction script (Sections~\ref{sec:lean} and~\ref{sec:repro}).
\item \textbf{Claim status.} A precise account of what is classical, what is cited, what is formalised and what may
  be new (Section~\ref{sec:status}).
\end{enumerate}

\section{Setting and main result}\label{sec:setting}

\subsection{Graphs, incidence matrices and currents}
A \emph{multigraph} $G$ has vertex set $V=\{0,1,\dots,n\}$ and edge set $\{1,\dots,q\}$. Each edge $i$ has an
orientation $\src(i)\to\tgt(i)$; loops and parallel edges are allowed. Every edge has unit resistance. The full
incidence matrix $M\in\R^{(n+1)\times q}$ has entries $M_{v,i}=[\src(i)=v]-[\tgt(i)=v]$, so a loop has a zero
column. The \emph{reduced} incidence matrix $A\in\R^{n\times q}$ deletes the row of the reference vertex~$0$.

If $G$ is connected, then $L=AA^{\T}$ (the reduced Laplacian) is invertible, and in fact positive definite. For
vertices $s,t$ let $b\in\R^n$ be the reduced vector with $+1$ at $s$ and $-1$ at $t$ (entries at vertex $0$
dropped). The unit current from $s$ to $t$ is
\[
I \;=\; A^{\T}L^{-1}b\in\R^q .
\]
With the potential $\eta=L^{-1}b$, extended by $\eta(0)=0$, Ohm's law reads $I_i=\eta(\src i)-\eta(\tgt i)$.
Kirchhoff's current law reads $AI=b$.

All entries of $A$ and $b$ are integers, so Cramer's rule gives
\[
I_e=\frac{N}{D},\qquad D=\det(AA^{\T})\in\Z_{>0},\qquad N=\bigl(A^{\T}\operatorname{adj}(AA^{\T})\,b\bigr)_e\in\Z .
\]
By the matrix-tree theorem, $D$ is the number $\tau(G)$ of spanning trees. By Kirchhoff's formula for currents,
$N$ is the number of spanning trees whose $s$--$t$ path uses $e$ forwards minus the number that use it backwards;
see \cite{Kirchhoff1847,DoyleSnell,LyonsPeres,Bollobas}. These identifications are classical and are \emph{not}
formalised here. The formal results are stated for $D$ and $N$ as defined above.

\subsection{Circulation networks and the tolerance contract}
A \emph{network} $K$ consists of a finite set of nodes and a finite set of arcs. Each arc $a$ has a tail, a head,
a capacity $\capac(a)\ge 0$ and a cost $c(a)\in\R$. All lower bounds are $0$. A circulation $x$ is \emph{feasible} if
$0\le x_a\le\capac(a)$ for all arcs and the net outflow is zero at \emph{every} node. Its cost is
$\cost(x)=\sum_a c(a)x_a$, and $\opt(K)=\inf\{\cost(x): x\text{ feasible}\}$.

\begin{definition}[Shared additive tolerance]\label{def:adm}
Let $K_A,K_B$ be two networks with a common set $O$ of observed arcs, and let $\varepsilon\in\R$. Write
$\Adm(\varepsilon)$ if there are feasible circulations $x^A$ of $K_A$ and $x^B$ of $K_B$ such that
\[
x^A|_O=x^B|_O,\qquad \cost_A(x^A)-\opt(K_A)\le\varepsilon,\qquad \cost_B(x^B)-\opt(K_B)\le\varepsilon .
\]
The \emph{threshold} of the pair is the least $\varepsilon$ with $\Adm(\varepsilon)$, if it exists.
\end{definition}

The two tolerances are additive and share a single $\varepsilon$. With separate tolerances one would obtain a
trade-off curve instead of a single number.

\subsection{The two networks}
Let $G^\ast$ be a multigraph on the same edge index set. Its vertices $\{0,\dots,k\}$ are called \emph{faces}, and
edge $i$ goes from $\src^\ast(i)$ to $\tgt^\ast(i)$. Let $B$ be its reduced incidence matrix.

In the planar case $G^\ast$ is the plane dual, \emph{with a consistent orientation}: every dual edge crosses its
primal edge in the same direction, say from the face on its left to the face on its right. In terms of the
rotation systems of Section~\ref{sec:cert}, dual edge $i$ goes from the face of the dart $(i,0)$ to the face of the
dart $(i,1)$. The orientation matters; see Remark~\ref{rem:orient}.

The \emph{shared topology} $\mathcal T(G,G^\ast,s,t)$ has the following nodes and arcs.
\begin{itemize}
\item \textbf{Nodes:} the vertices of $G$, the faces of $G^\ast$, and two new nodes $u_i,v_i$ for each edge $i$.
\item \textbf{Observed arcs:} $o_i\colon u_i\to v_i$ for every edge $i$.
\item \textbf{Primal connectors:} $\src(i)\to u_i$ and $v_i\to\tgt(i)$.
\item \textbf{Dual connectors:} $\src^\ast(i)\to u_i$ and $v_i\to\tgt^\ast(i)$.
\item \textbf{Return arc:} one arc $r\colon t\to s$.
\end{itemize}
The two networks use this topology with the following capacities and costs:
\begin{center}
\begin{tabular}{lccccl}
\toprule
 & observed $o_i$ & primal conn. & dual conn. & return $r$ & costs\\
\midrule
network $\mathcal A$ & $1$ & $1$ & $0$ & $1$ & $c(r)=-1$, all others $0$\\
network $\mathcal B_e$ & $1$ & $0$ & $1$ & $0$ & $c(o_e)=+1$, all others $0$\\
\bottomrule
\end{tabular}
\end{center}
The auxiliary topology need not be planar, and its size is linear in $q$.

\paragraph{Orientation.} Let $\sigma_i=-1$ if $I_i<0$ and $\sigma_i=+1$ otherwise. Write $G^\sigma$ and
$G^{\ast\sigma}$ for the multigraphs obtained by reversing every edge $i$ with $\sigma_i=-1$, \emph{both} in $G$
and in $G^\ast$. Reversing an edge negates its current, so in $G^\sigma$ all currents are nonnegative and the
current on $e$ is $|I_e|$.

\begin{theorem}[Exact threshold]\label{thm:main}
Let $G$ be a connected planar multigraph with plane dual $G^\ast$ (for some plane embedding), oriented
consistently as above. Let $s,t$ be vertices
and let $e$ be an edge with $I_e=I_e(s,t)\neq 0$. Let $\mathcal A$ and $\mathcal B_e$ be the networks above, built
on $\mathcal T(G^\sigma,G^{\ast\sigma},s,t)$. Then the threshold of $(\mathcal A,\mathcal B_e)$ exists and equals
\[
\Delta^*(G;s,t,e)=\frac{|I_e|}{1+|I_e|}=\frac{|N|}{D+|N|}.
\]
\end{theorem}

The only property of $G^\ast$ that the proof uses is
\begin{equation}\label{eq:H2}
\{z\in\R^q : Bz=0\}=\{A^{\T}y : y\in\R^n\}\qquad(\ker B=\row A). \tag{H2}
\end{equation}
Theorem~\ref{thm:main} therefore holds verbatim for any connected $G$ and any multigraph $G^\ast$ on the same edge
set that satisfies \eqref{eq:H2}. This general form is the one proved first in Lean (\code{theorem1\_networks}).
Property \eqref{eq:H2} is preserved when the same set of edges is reversed in $G$ and in $G^\ast$.

\begin{remark}\label{rem:orient}
Property \eqref{eq:H2}, and with it Theorem~\ref{thm:main}, depends on the orientation of the dual edges. Take
the digon, two parallel edges $0\to1$, with $s=0$ and $t=1$, so that $I=(\tfrac12,\tfrac12)$. With the consistent
orientation the two dual edges are antiparallel, $\ker B=\row A$, and the threshold is $1/3$. If both dual edges
point from the inner to the outer face, then $\ker B$ is spanned by $(1,-1)$, the only common observation is
$z=0$, and the threshold is $1$.
\end{remark}

\begin{corollary}\label{cor:theta}
Let $\mu(G)$ be the least nonzero value of $|I_e(s,t)|$ over all $s\neq t$ and all edges $e$. The least
threshold over all admissible $(s,t,e)$ is $\Theta(G)=\mu(G)/(1+\mu(G))$, and $\Theta(G)\ge 1/(D+1)$. Equality
holds if and only if some nonzero numerator has $|N|=1$.
\end{corollary}
\begin{proof}
The function $x\mapsto x/(1+x)$ is increasing. Moreover $|I_e|=|N|/D\ge 1/D$ whenever $N\neq 0$.
\end{proof}

\begin{example}\label{ex:fan}
For the fan with a hub joined to a path $p_1\cdots p_k$ ($q=2k-1$ edges), feeding $p_1$ against the hub gives
current $1/F_{2k}$ on the last path edge, where $F_m$ are the Fibonacci numbers. Here $D=F_{2k}$ and $|N|=1$, so
$\Delta^*=1/(F_{2k}+1)$. Our independent check in Section~\ref{sec:repro} recomputes $D=3,8,21,55$ and the
minimal thresholds $1/4,1/9,1/22,1/56$ for $k=2,3,4,5$. For $K_4$ with one doubled edge, $D=24$ and the minimal
threshold is $1/25$.
\end{example}

\section{Proof of Theorem~\ref{thm:main}}\label{sec:proof}

Throughout, $G$ is connected, $G^\ast$ satisfies \eqref{eq:H2}, and we work with $G^\sigma$ and $G^{\ast\sigma}$.
For readability we keep writing $A$, $B$, $I$ for their matrices and currents. We write $a:=I_e=|I_e|>0$.

\begin{lemma}[Unit current bounds]\label{lem:01}
In $G^\sigma$ we have $0\le I_i\le 1$ for every edge $i$.
\end{lemma}
\begin{proof}
Nonnegativity holds by the choice of $\sigma$. For the upper bound, fix $c\in\R$ and put
$\chi(w)=[\eta(w)\ge c]$. Each edge satisfies $I_i\,(\chi(\src i)-\chi(\tgt i))\ge 0$, because
$I_i=\eta(\src i)-\eta(\tgt i)$ and $\chi$ is a nondecreasing function of $\eta$. Summing over $i$ and using
Kirchhoff's law with the full incidence matrix ($MI$ equals $+1$ at $s$, $-1$ at $t$, and $0$ elsewhere) gives
\[
\sum_i I_i\,(\chi(\src i)-\chi(\tgt i))=\sum_w\chi(w)\,(MI)_w=\chi(s)-\chi(t)\le 1 .
\]
Every summand is nonnegative, so each is at most~$1$. For an edge $i_0$ with $I_{i_0}>0$, take
$c=\eta(\src i_0)$. Then $\chi(\src i_0)=1$ and $\chi(\tgt i_0)=0$, hence $I_{i_0}\le 1$.
\end{proof}

The bound is used below to show that $(z,f)=(I,1)$ is feasible in $\mathcal A$. In the graph setting it always
holds, but the abstract assembly step needs it as a hypothesis. For $A=[1]$ (one edge) and the two-unit feed
$b=[2]$, which is not a unit $s$--$t$ feed, one gets $I=2$. The optimum of $\mathcal A$ then becomes $-1/2$ and the
gap formula below fails. This example is formalised in \code{Necessity.lean}.

\begin{lemma}[Projections]\label{lem:proj}
For network $\mathcal A$, the map $x\mapsto(x|_O,x_r)$ sends the feasible circulations onto
\[
P_{\mathcal A}=\{(z,f)\in\R^q\times\R:\ Az=fb,\ 0\le z\le 1,\ 0\le f\le 1\},
\]
and $\cost_{\mathcal A}(x)=-x_r$. For network $\mathcal B_e$, the map $x\mapsto x|_O$ sends the feasible
circulations onto $P_{\mathcal B}=\{z: Bz=0,\ 0\le z\le 1\}$, and $\cost_{\mathcal B_e}(x)=x_{o_e}$.
\end{lemma}
\begin{proof}
At $u_i$ the only arcs are the observed arc and the incoming connectors. Connectors of capacity $0$ carry nothing.
Balance at $u_i$ therefore forces the active incoming connector to carry $z_i:=x_{o_i}$, and in the same way at
$v_i$ for the outgoing connector.

In $\mathcal A$, let $f:=x_r$. The net outflow at a vertex $w$ of $G$ is $(Mz)_w-f\,([w=s]-[w=t])$. All these
vanish if and only if $Az=fb$: the row of the reference vertex is redundant, because the columns of $M$ and the
full right-hand side both sum to zero. Face nodes carry no flow. In $\mathcal B_e$ the vertex nodes carry no flow,
and balance at the faces is $Bz=0$ by the same argument. Loops of either graph need no special treatment: both
connectors of a loop meet the same node and cancel, matching its zero incidence column.

Conversely, every $(z,f)\in P_{\mathcal A}$ (respectively $z\in P_{\mathcal B}$) extends uniquely to a feasible
circulation, by putting $z_i$ on the active connectors of $i$ and $f$ on $r$. The cost formulas are immediate.
\end{proof}

\begin{lemma}[Optima]\label{lem:opt}
$\opt(\mathcal A)=-1$ and $\opt(\mathcal B_e)=0$.
\end{lemma}
\begin{proof}
On $P_{\mathcal A}$ the cost is $-f\ge -1$. The value $-1$ is attained at $(I,1)$, which lies in $P_{\mathcal A}$
by $AI=b$ and Lemma~\ref{lem:01}. On $P_{\mathcal B}$ the cost is $z_e\ge 0$, attained at $z=0$.
\end{proof}

\begin{lemma}[Common observations]\label{lem:common}
If $(z,f)\in P_{\mathcal A}$ and $z\in P_{\mathcal B}$, then $z=fI$. Conversely, for every $f\in[0,1]$ we have
$(fI,f)\in P_{\mathcal A}$ and $fI\in P_{\mathcal B}$.
\end{lemma}
\begin{proof}
If $Bz=0$, then $z=A^{\T}y$ for some $y$ by \eqref{eq:H2}. Then $AA^{\T}y=fb$, so $y=fL^{-1}b$ and $z=fI$.

Conversely, $A(fI)=fb$. Also $B(fI)=f\,BA^{\T}L^{-1}b=0$, since $BA^{\T}=0$ by the inclusion
$\row A\subseteq\ker B$ in \eqref{eq:H2}. The bounds follow from Lemma~\ref{lem:01} and $0\le f\le 1$.
\end{proof}

\begin{proof}[Proof of Theorem~\ref{thm:main}]
By Lemmas~\ref{lem:proj}--\ref{lem:common}, $\Adm(\varepsilon)$ holds if and only if some $f\in[0,1]$ satisfies
\[
1-f\le\varepsilon\qquad\text{and}\qquad af\le\varepsilon .
\]
The first inequality bounds the gap of $\mathcal A$ at $(fI,f)$, which is $-f-(-1)$. The second bounds the gap of
$\mathcal B_e$ at $fI$, which is $af-0$. The function $\max(1-f,af)$ is minimised on $[0,1]$ at $f^*=1/(1+a)$,
where both terms equal $a/(1+a)$. Hence $\Adm(\varepsilon)$ holds exactly when $\varepsilon\ge a/(1+a)$, and the
threshold is $a/(1+a)$.

Finally $a=|N|/D$ with $D>0$, so $a/(1+a)=|N|/(D+|N|)$. It remains to note that the consistently oriented plane dual satisfies
\eqref{eq:H2}: its cycle space is the cut space of $G$. Section~\ref{sec:cert} derives this from a rotation-system
certificate.
\end{proof}

\begin{remark}
The exact optima of the two networks have disjoint observations. The optimum of $\mathcal A$ forces $f=1$ and
$z=I$, so $z_e=a>0$, while the optimum of $\mathcal B_e$ forces $z_e=0$. This is why the threshold is positive.
\end{remark}

\section{Rotation-system certificates}\label{sec:cert}

Planarity is encoded by a finite, decidable certificate. A \emph{dart} is a pair $(i,\beta)$ with $\beta\in\{0,1\}$:
the dart $(i,0)$ sits at $\src(i)$ and the dart $(i,1)$ sits at $\tgt(i)$. The involution $\alpha$ swaps the two
darts of an edge.

\begin{definition}[Genus-$0$ rotation-system certificate]\label{def:rotcert}
A certificate $C$ with $k+1$ faces for a multigraph $G$ (the parameter \code{f} in Lean) consists of a permutation $\rot$ of the darts and a labelling
$\face$ of the darts by $\{0,\dots,k\}$. They must satisfy six conditions:
\begin{enumerate}\renewcommand{\labelenumi}{(\roman{enumi})}
\item $\rot$ fixes the vertex of every dart;
\item $\face$ is invariant under $\varphi=\rot\circ\alpha$;
\item the darts at each vertex form a single $\rot$-orbit;
\item equal face labels occur only within one $\varphi$-orbit;
\item every label occurs;
\item $(n+1)-q+(k+1)=2$.
\end{enumerate}
The \emph{dual} $C^\ast$ has the faces as vertices; edge $i$ goes from $\face(i,0)$ to $\face(i,1)$.
\end{definition}

Conditions (ii), (iv) and (v) say that the labels are exactly the $\varphi$-orbits. Conditions (i) and (iii) say
that the $\rot$-orbits are exactly the vertices that carry darts. Condition (vi) is Euler's equation for genus~$0$.

\begin{proposition}\label{prop:H2}
Let $G$ be connected with a certificate $C$, and let $B$ be the reduced incidence matrix of $C^\ast$. Then:
\begin{enumerate}\renewcommand{\labelenumi}{(\alph{enumi})}
\item $BA^{\T}=0$;
\item $C^\ast$ is connected;
\item $\ker B=\row A$, which is \eqref{eq:H2}.
\end{enumerate}
\end{proposition}
\begin{proof}
(a) Fix a vertex $v$ and a face $g$. Expanding the incidence entries over darts gives
\[
\sum_i M_{v,i}M^\ast_{g,i}=\sum_{d\text{ at }v}\bigl([\face d=g]-[\face\alpha d=g]\bigr).
\]
Each dart at $v$ contributes once; a loop contributes two darts that cancel. Condition (ii) gives
$\face(\alpha d)=\face(\rot d)$. Since $\rot$ permutes the darts at $v$, the two sums agree and the total is $0$.
Passing to submatrices gives $BA^{\T}=0$.

(b) The faces $\face d$ and $\face\alpha d$ are equal or joined by the dual edge of $d$, and
$\face(\rot d)=\face(\alpha d)$. So consecutive faces along any sequence of $\alpha$- and $\rot$-steps are equal or
adjacent. Since $G$ is connected, every walk in $G$ lifts, by (iii), to such a dart sequence. Hence
$\langle\rot,\alpha\rangle$ is transitive on darts. By (v) every label is the face of some dart, so any two
vertices of $C^\ast$ are joined by a walk.

(c) Connectivity of $G$ makes $AA^{\T}$ invertible, so $A^{\T}$ is injective and $\dim\row A=n$. Connectivity of
$C^\ast$ (with $k+1$ vertices) makes $BB^{\T}$ invertible, so $B$ is onto $\R^k$ and $\dim\ker B=q-k$. Euler (vi)
gives $n=q-k$. By (a), $\row A\subseteq\ker B$; equal dimensions give equality.
\end{proof}

Condition (iv) is not used in this proof. It is part of what the certificate means: it guarantees that $k+1$ is
the number of faces.

\begin{proposition}[cited, not formalised]\label{prop:cited}
A connected finite multigraph with at least one edge is planar if and only if it admits a genus-$0$ rotation-system
certificate. For a plane embedding, the clockwise rotation yields a certificate whose dual is the plane dual,
with the consistent orientation of Section~\ref{sec:setting}.
\end{proposition}

This is the Heffter--Edmonds--Ringel rotation principle together with the classification of surfaces and Euler's
formula; see \cite{MoharThomassen,GrossTucker}. In one direction, a plane embedding gives a rotation whose
$\varphi$-orbits are the face boundaries, and Euler's formula gives (vi). In the other direction, every rotation
system is a cellular embedding in an orientable surface with Euler characteristic $(n+1)-q+(k+1)$, and
characteristic $2$ means the sphere. The one-vertex graph without edges has no certificate in this encoding, because
labels come from darts; it is irrelevant here. In Lean, ``planar'' \emph{means} ``has a certificate'', and
Theorem~\ref{thm:main} covers all connected planar multigraphs only through Proposition~\ref{prop:cited}.

\begin{example}
The Lean file \code{Final.lean} contains a certificate for a triangle with a pendant edge (4~vertices,
4~edges, 2~faces). Its conditions are checked by \code{decide}, and Euler's equation by \code{norm\_num}. Its dual has three parallel edges and a
loop, the loop being the dual of the bridge. Instantiating the main theorem with $s=0$, $t=3$ and $e$ the bridge
gives threshold $1/2$.
\end{example}

\section{The Lean formalisation}\label{sec:lean}

\subsection{Main statement}
The integrated theorem is \code{PlanarCore.theorem1\_certified} in \code{PlanarCore/Final.lean}:
\begin{lstlisting}
theorem theorem1_certified {n q f : ℕ} (G : MGraph n q) (hG : G.Connected) (C : RotCert G f)
    (s t : Fin (n+1)) (e : Fin q) (he : current G.redInc (redB s t) e ≠ 0) :
    0 < G.detL ∧
    current G.redInc (redB s t) e = (G.numer s t e : ℝ) / (G.detL : ℝ) ∧
    IsLeast {ε | AdmissibleTransfer
        (netA (G.oriented s t) (C.dual.reorient (fun i => current G.redInc (redB s t) i < 0)) s t)
        (netB (G.oriented s t) (C.dual.reorient (fun i => current G.redInc (redB s t) i < 0)) s t e)
        ε}
      ((|G.numer s t e| : ℝ) / ((G.detL : ℝ) + |G.numer s t e|))
\end{lstlisting}
Its only hypotheses are the following. \code{hG} is connectivity, every pair of vertices being joined by a walk.
\code{C} is the certificate of Definition~\ref{def:rotcert}. \code{he} says that the observed edge carries
current. The quantities \code{detL} and \code{numer} are the integers $D$ and $N$. The networks \code{netA} and
\code{netB} are the explicit networks of Section~\ref{sec:setting}. \code{AdmissibleTransfer} is
Definition~\ref{def:adm}, with optima computed as infima over feasible circulations.

\subsection{Modules}
\begin{center}
\small
\begin{tabularx}{\linewidth}{>{\ttfamily}l r X}
\toprule
\textrm{module} & lines & content\\
\midrule
Graph & 241 & multigraph model; incidence; currents; $AA^{\T}$ invertible for connected $G$; $|I_i|\le 1$\\
Threshold & 123 & minimum of two linear gaps; arithmetic of the threshold\\
Projection, Assembly & 43, 112 & abstract projection algebra; assembly with \eqref{eq:H2} as hypothesis\\
Network & 378 & explicit networks \code{netA}, \code{netB}; projections, costs, optima; \code{theorem1\_networks}\\
Certificate & 226 & \code{RotCert}; dual; $BA^{\T}=0$; connectivity of the dual; \eqref{eq:H2}\\
Kirchhoff & 96 & $D>0$; $I_e=N/D$ by Cramer's rule\\
Final & 93 & \code{theorem1\_certified}; decided example\\
Theorem1, Sanity & 103, 76 & intermediate integration, including $0\le I\le1$ after reorientation; non-vacuity\\
Necessity & 43 & counterexample showing $0\le I\le1$ is needed\\
Fekete & 164 & auxiliary: abstract Fekete step, substitution circuit algebra, AM--GM inequality (Appendix~\ref{app:growth})\\
\bottomrule
\end{tabularx}
\end{center}

\subsection{Design choices and remaining conventions}
Vertices are \code{Fin (n+1)}, with vertex $0$ as reference, and edges are \code{Fin q}. Darts are
\code{Fin q × Bool}. Bridges (a zero column of $B$, that is a dual loop), loops (a zero column of $A$) and parallel
edges need no special case. Strict positivity $D>0$ comes from positive definiteness of the Gram matrix.

One convention is not separately proved. The networks are built from $G^\sigma$ and $(C^\ast)^\sigma$, where $C$ is
a certificate for $G$. That $(C^\ast)^\sigma$ is the dual of a certificate for $G^\sigma$ is immediate from the
definitions: swap the two darts of each reversed edge. This fact is not part of the checked statement. The
statement as formalised concerns the networks built from $G^\sigma$ and $(C^\ast)^\sigma$, and it does not depend
on this fact.

\section{Reproduction and checks}\label{sec:repro}
\sloppy

\paragraph{Versions.} Lean \code{leanprover/lean4:v4.37.0-rc1}; Mathlib commit\\
\code{53dad9ed98be8416b6ca1111e55a03de4303231f}, pinned in \code{lakefile.toml} and \code{lake-manifest.json}.

\paragraph{Commands.} \code{./scripts/check.sh} runs five steps. It fetches the prebuilt Mathlib cache
(\code{lake exe cache get}) and builds the library (\code{lake build}). It then searches the sources for
\code{sorry}, \code{admit} and \code{axiom}, prints the axioms of the main declarations
(\code{scripts/PrintAxioms.lean}), and replays every module with \code{leanchecker}.

\paragraph{Recorded results (development environment).} The checks gave the following results.
\begin{itemize}
\item \code{lake build} after deleting \code{.lake/build}: \emph{Build completed successfully (18329 jobs)}, exit
  code~$0$. The only warning is an advisory from Mathlib.
\item The search for \code{sorry}, \code{admit} and \code{axiom} found nothing.
\item \code{\#print axioms} gives \code{[propext, Classical.choice, Quot.sound]} for \code{theorem1\_certified},
  \code{RotCert.H2}, \code{RotCert.orth}, \code{RotCert.dual\_mul\_transpose}, \code{detL\_pos},
  \code{current\_eq\_numer\_div}, \code{threshold\_cramer}, \code{triCert}, \code{triG\_certified\_threshold},
  \code{theorem1\_networks}, \code{projA}, \code{projB}, \code{H1\_of\_connected}, \code{current\_abs\_le\_one} and
  \code{T3\_fekete}. \code{RotCert.dual\_connected} uses only \code{[propext, Quot.sound]}.
\item \code{leanchecker} replays all 12 modules with exit code~$0$.
\end{itemize}
The clean-room run of the published package is recorded in \code{VERIFICATION.md}. The package was copied into an
empty directory, Mathlib was cloned and its cache downloaded afresh, and the script was run unchanged. The run used
the same machine and the already installed Lean toolchain binary. The script prints the axioms without checking
them; they were inspected by hand.

\paragraph{Separate numerical cross-check (not part of the proof).} We computed the threshold separately by
linear programming, with a separate implementation (\code{crosscheck/lp\_crosscheck.py}, written by the same AI
agents as the rest of the project). A plane embedding gives the dual. We built $\mathcal A$ and $\mathcal B_e$ as above, computed
both optima and the least shared tolerance with the HiGHS solver, and compared the result with $|I_e|/(1+|I_e|)$
computed in exact rational arithmetic. The comparison covered all $1556$ admissible triples $(s,t,e)$ on eight
planar multigraphs:
\begin{itemize}
\item the triangle with a pendant edge;
\item $K_4$ with one doubled edge;
\item the fans with $k=2,\dots,5$;
\item the octahedron;
\item the cube.
\end{itemize}
None of these graphs has loops. The largest deviation was $1.1\times10^{-15}$, both optima were $-1$ and $0$ in every case, and the minimal
thresholds agree with Example~\ref{ex:fan}.

\section{Related work}\label{sec:related}

\paragraph{Currents and spanning trees.} The formula for $I_e$ as a ratio of spanning-tree counts goes back to
Kirchhoff \cite{Kirchhoff1847}; modern accounts are \cite{DoyleSnell,Bollobas,LyonsPeres}. The duality
$\ker B=\row A$ between the cycle space of the dual and the cut space of a plane graph is classical. Bu and Pan
\cite{BuPan} study square matrices obtained by concatenating two incidence-type blocks. They give a
linear-algebraic planarity criterion through the number of spanning trees: an excess, defined from the
spanning-tree count and a maximum determinant, is $0$ for planar graphs and at least $18$ otherwise. We could not
access the full text during this work and describe it from its abstract and public summaries. For a planar graph
the stacked matrix $[A;B]$ of primal and dual reduced incidences satisfies $|\det[A;B]|=\tau(G)$. This follows
from $BA^{\T}=0$ and the matrix-tree theorem, since $\det([A;B][A;B]^{\T})=\det(AA^{\T})\det(BB^{\T})=
\tau(G)\tau(G^\ast)=\tau(G)^2$. Our threshold does not use this determinant, but it lives in the same primal--dual
setting.

\paragraph{Inverse problems for resistances.} Chan, Kontorovich and Pak \cite{CKP25} determine, up to a constant
factor, the smallest simple planar graph with a prescribed effective resistance $\tau(G/e)/\tau(G)$. Their earlier
work \cite{CKP24} treats spanning-tree counts. The effective resistance of an edge is the current through that edge
when its own endpoints are the terminals. In a simple graph it is at least
$\tfrac12(1/\deg x+1/\deg y)$ for an edge $xy$ \cite[Lemma~3.4]{CKP25}, hence of order at least $1/|V|$. In contrast, the quantity $\mu(G)$ of Corollary~\ref{cor:theta} allows arbitrary terminals
and can be exponentially small; for fans it equals $1/F_{2k}$.

\paragraph{Squared rectangles.} For terminals on a common face, the currents of a planar network of unit resistors
are the side lengths of a tiling of a rectangle by squares, one square for each edge that carries current, and the
total current is the width of the rectangle \cite{BSST}. Hence, for co-facial terminals, $1/\mu$ is the ratio of
the width of the rectangle to its smallest square, and the quantity $\Lambda^{\mathrm{face}}_q$ of the Outlook and
of Appendix~\ref{app:growth} is the largest such ratio, over both orientations, for tilings by at most $q$ squares.
This classical literature contains the known bounds. Conway \cite{Conway} and Kenyon \cite{Kenyon} give
logarithmic lower bounds on the number of squares, that is, exponential upper bounds with base $2$ on such
ratios. Kenyon's Fibonacci tiling realises the golden ratio, and he asks for the best constants. Fomin \cite{Fomin}
proves $\tau(G)\le1.8637^q$ for planar multigraphs with $q$ edges and deduces that an integer rectangle containing
a unit square and with longest side at least $1.8637^q$ needs at least $q$ squares; in our notation,
$\Lambda^{\mathrm{face}}_q\le1.8637^q$. Keleti, Lacina, Liu, Liu and Tuir\'an Rangel \cite{Keleti} bound the
minimal integer scaling of a tiling by $n$ squares by $4^n$ and cite the Fibonacci tiling as the known witness from
below. Version~1.0 of this paper omitted this literature.

\paragraph{Equal-flow problems.} Network problems with side constraints that force flows on designated arcs to be
equal have been studied at least since \cite{AliKenningtonShetty}; see \cite{AhujaOrlinSechiZuddas} for the simple
equal-flow problem. Those works optimise one network under equality constraints. Our question is different: we
compare two networks with \emph{different} costs and capacities and ask for the least shared optimality gap at a
common observation.

\paragraph{Formalised planarity.} The four-colour theorem was formalised in Coq using hypermaps
\cite{Gonthier2008}. Isabelle/HOL has a formalisation of planarity through combinatorial maps, with verified checkers
for planarity certificates \cite{NoschinskiAFP}. Certifying graph algorithms with planarity certificates are
discussed in \cite{AbdulazizMehlhornNipkow}. Our certificate (Definition~\ref{def:rotcert}) follows the same
combinatorial idea in Lean~4 with Mathlib \cite{mathlib,lean4}. We use it to derive the linear-algebraic duality
\eqref{eq:H2} needed for electrical networks.

\section{Status of the claims}\label{sec:status}

\begin{center}
\small
\begin{tabularx}{\linewidth}{l X}
\toprule
classical & Kirchhoff's laws and current formula; the matrix-tree theorem; planar duality and Euler's formula; the
rotation principle; Fekete's lemma\\
prior work & planarity criterion via spanning-tree numbers and concatenated incidence matrices \cite{BuPan}; inverse
problems for effective resistance and spanning-tree counts \cite{CKP25,CKP24}; squared rectangles and bounds on
their size \cite{BSST,Conway,Kenyon,Fomin,Keleti}\\
cited, not formalised & planarity $\Leftrightarrow$ genus-$0$ rotation-system certificate (Proposition~\ref{prop:cited});
$D=\tau(G)$ and $|N|$ as a signed tree count (classical)\\
proved informally, not formalised & that $(C^\ast)^\sigma$ is the dual of a certificate for $G^\sigma$
(Section~\ref{sec:lean}); the minimum over $(s,t,e)$ (Corollary~\ref{cor:theta}); the fan formula for general $k$
(Example~\ref{ex:fan}); the substitution lemma (Appendix~\ref{app:growth})\\
formalised (Lean) & Theorem~\ref{thm:main} for certified graphs: $D>0$, $I_e=N/D$, $0\le I\le1$, the networks,
their projections and optima, $BA^{\T}=0$, connectivity of the dual, \eqref{eq:H2}, the exact threshold; a decided
example; the abstract Fekete step\\
possibly new (no priority claim) & the two-network common-observation threshold as a packaged statement with its
realisation by circulations; the formal derivation of \eqref{eq:H2} from genus-$0$ rotation certificates in
Lean/Mathlib\\
\bottomrule
\end{tabularx}
\end{center}

Our literature search was limited. For version 1.0 it covered web searches and the full texts of
\cite{CKP25,CKP24}; for \cite{BuPan} only the abstract and public summaries were available. For version 1.1 we
added the squared-rectangle literature \cite{BSST,Conway,Kenyon,Fomin,Keleti}. We found no earlier statement of Theorem~\ref{thm:main}. We regard the
result as mathematically light, a corollary of classical facts, and we do not claim priority. We would be grateful
for references to earlier occurrences.

\section{Outlook}

The quantity $\mu(G)$ suggests an extremal problem. Let $\Lambda_q$ be the largest value of $1/\mu(G)$ over connected
planar multigraphs with $q$ edges, and let $\Lambda^{\mathrm{face}}_q$ be the same with $s,t$ on a common face.
Appendix~\ref{app:growth} records an elementary supermultiplicativity property and the resulting existence of
$\lim_q(\Lambda^{\mathrm{face}}_q)^{1/q}$. The value of this growth constant, and whether $\Lambda_q$ grows at the
same exponential rate, are open as far as we know. They are the subject of ongoing work. For co-facial terminals
this is the classical problem of squared rectangles; see Section~\ref{sec:related}.

\section*{Author contributions and use of AI}
This work was carried out with substantial assistance from AI systems working under the author's direction. The
SAPIO MORPHEUS project uses models from OpenAI and Anthropic.
\begin{itemize}
\item The mathematical write-up behind Sections~\ref{sec:setting}--\ref{sec:proof} was drafted by one AI agent of
  the project.
\item The Lean formalisation, the literature comparison, the numerical cross-check and this manuscript were
  produced by AI agents based on Anthropic's Claude (Claude Code).
\item The author selected the problem, set the acceptance criteria, directed the work and made the publication
  decisions. The author takes responsibility for the content.
\end{itemize}
Every formal claim was checked by the Lean kernel. The informal parts were reviewed for consistency with the
formal statements. In line with common publication policies, AI systems are not listed as authors.

\appendix
\section{Auxiliary: growth of the inverse minimal current}\label{app:growth}

This appendix is informal. Only the abstract Fekete step, the circuit algebra of the substitution lemma and an
AM--GM inequality are formalised, in \code{Fekete.lean}, with the graph-theoretic inputs as hypotheses.

For $q\ge1$ let $\Lambda^{\mathrm{face}}_q$ be the largest value of $1/|I_e(s,t)|$ over connected planar
multigraphs with $q$ edges, terminals $s\neq t$ on a common face of some plane embedding, and edges $e$ with
$I_e\neq0$. Equivalently, $s,t$ are co-facial if and only if $G+st$ is planar.

\begin{lemma}[Substitution]
Let $G$ have an observed edge $uv$ with $I_{uv}\neq0$. Let $H$ be connected with distinct ports $a,b$ on its outer
face and an edge $e'$ with nonzero current $i_H$ under the unit $a$--$b$ feed. Replace $uv$ by a unit resistor from
$u$ to $a$ in series with $H$, identifying $b$ with $v$. The result is planar and has $q_G+q_H$ edges, and the current
on $e'$ is $J\,i_H$ with $0<|J|\le|I_{uv}|$. Co-faciality of $s,t$ is preserved.
\end{lemma}
\begin{proof}[Sketch]
If $uv$ is not a bridge, let $R_0$ be the resistance of $G-uv$ between $u$ and $v$, and let $V_0$ be the open-circuit
voltage. The inserted branch has resistance $R=1+R_H(a,b)\ge1$, so $J=V_0/(R_0+R)$, while $I_{uv}=V_0/(R_0+1)$.
If $uv$ is a bridge, both currents equal $\pm1$. Linearity gives the factor $i_H$. The new branch can be drawn
inside a thin neighbourhood of the removed edge.
\end{proof}

Hence $\Lambda^{\mathrm{face}}_{p+q}\ge\Lambda^{\mathrm{face}}_p\Lambda^{\mathrm{face}}_q$. Moreover
$1\le\Lambda^{\mathrm{face}}_q\le T_q$, where $T_q$ is the largest number of spanning trees of a planar
multigraph with $q$ edges: for a maximising graph $G$ we have $|I_e|=|N|/D\ge1/D$ and $D=\tau(G)\le T_q$. The
number $T_q$ grows at most exponentially. The inequality $(x/j)^j\le\exp(x/\mathrm e)$, which is formalised, gives
$T_q\le\exp(2q/\mathrm e)$ once combined with the matrix-tree theorem and $\operatorname{tr}L\le2q$; that
combination is not formalised. Fekete's lemma \cite{Fekete} then gives the existence of
$\lim_q(\Lambda^{\mathrm{face}}_q)^{1/q}=\sup_q(\Lambda^{\mathrm{face}}_q)^{1/q}$ (\code{T3\_fekete}). The fans of
Example~\ref{ex:fan} show that this limit is at least the golden ratio $(1+\sqrt5)/2$. Its exact value is open as
far as we know. In the language of squared rectangles (Section~\ref{sec:related}) the co-facial quantity is the
classical ratio of the width of a squared rectangle to its smallest square. The base~$2$ and Fomin's sharper base
$1.8637$ \cite{Fomin} are known upper bounds, and the golden-ratio tiling is classical \cite{Kenyon,Keleti}. We
found no earlier statement that the limit exists.

\bibliographystyle{amsplain}
\begin{thebibliography}{99}
\bibitem{AbdulazizMehlhornNipkow} M.~Abdulaziz, K.~Mehlhorn, T.~Nipkow, \emph{Trustworthy graph algorithms},
  44th Int.\ Symp.\ Mathematical Foundations of Computer Science (MFCS 2019), LIPIcs 138, 1:1--1:22;
  arXiv:1907.04065.
\bibitem{AhujaOrlinSechiZuddas} R.~K.~Ahuja, J.~B.~Orlin, G.~M.~Sechi, P.~Zuddas, \emph{Algorithms for the simple
  equal flow problem}, Management Science \textbf{45} (1999), 1440--1455.
\bibitem{AliKenningtonShetty} A.~I.~Ali, J.~Kennington, B.~Shetty, \emph{The equal flow problem}, European J.\
  Oper.\ Res.\ \textbf{36} (1988), 107--115.
\bibitem{Bollobas} B.~Bollob\'as, \emph{Modern Graph Theory}, Graduate Texts in Mathematics 184, Springer, 1998.
\bibitem{BSST} R.~L.~Brooks, C.~A.~B.~Smith, A.~H.~Stone, W.~T.~Tutte, \emph{The dissection of rectangles into
  squares}, Duke Math.\ J.\ \textbf{7} (1940), 312--340.
\bibitem{BuPan} A.~Bu, Y.~Pan, \emph{Planarity via spanning tree number: a linear-algebraic criterion}, SIAM J.\
  Discrete Math.\ \textbf{39} (2025), 728--751; doi:10.1137/24M1660395.
\bibitem{CKP24} S.~H.~Chan, A.~Kontorovich, I.~Pak, \emph{Spanning trees and continued fractions},
  arXiv:2411.18782 (2024).
\bibitem{CKP25} S.~H.~Chan, A.~Kontorovich, I.~Pak, \emph{Effective resistance in planar graphs and continued
  fractions}, arXiv:2505.19168v1 (2025).
\bibitem{Conway} J.~H.~Conway, \emph{Mrs Perkins's quilt}, Proc.\ Cambridge Philos.\ Soc.\ \textbf{60} (1964),
  363--368.
\bibitem{lean4} L.~de~Moura, S.~Ullrich, \emph{The Lean 4 theorem prover and programming language}, Automated
  Deduction -- CADE 28, LNCS 12699, Springer, 2021, 625--635.
\bibitem{DoyleSnell} P.~G.~Doyle, J.~L.~Snell, \emph{Random Walks and Electric Networks}, Carus Math.\ Monographs
  22, MAA, 1984.
\bibitem{Fekete} M.~Fekete, \emph{\"Uber die Verteilung der Wurzeln bei gewissen algebraischen Gleichungen mit
  ganzzahligen Koeffizienten}, Math.\ Z.\ \textbf{17} (1923), 228--249.
\bibitem{Fomin} D.~Fomin, \emph{Elementary proof for the bounds of the complexity of a planar multigraph and the
  size of a prime rectangular squaring}, arXiv:2103.10523 (2021).
\bibitem{Gonthier2008} G.~Gonthier, \emph{Formal proof---the four-color theorem}, Notices Amer.\ Math.\ Soc.\
  \textbf{55} (2008), 1382--1393.
\bibitem{GrossTucker} J.~L.~Gross, T.~W.~Tucker, \emph{Topological Graph Theory}, Wiley, 1987.
\bibitem{Keleti} T.~Keleti, S.~Lacina, C.~Liu, M.~Liu, J.~R.~Tuir\'an Rangel, \emph{Tiling of rectangles with
  squares and related problems via Diophantine approximation}, Discrete Math.\ \textbf{346} (2023), 113442.
\bibitem{Kenyon} R.~Kenyon, \emph{Tiling a rectangle with the fewest squares}, J.\ Combin.\ Theory Ser.\ A
  \textbf{76} (1996), 272--291.
\bibitem{Kirchhoff1847} G.~Kirchhoff, \emph{\"Uber die Aufl\"osung der Gleichungen, auf welche man bei der
  Untersuchung der linearen Vertheilung galvanischer Str\"ome gef\"uhrt wird}, Ann.\ Phys.\ Chem.\ \textbf{72}
  (1847), 497--508.
\bibitem{LyonsPeres} R.~Lyons, Y.~Peres, \emph{Probability on Trees and Networks}, Cambridge University Press,
  2016.
\bibitem{mathlib} The mathlib Community, \emph{The Lean mathematical library}, Proc.\ 9th ACM SIGPLAN Int.\ Conf.\
  Certified Programs and Proofs (CPP 2020), 367--381.
\bibitem{MoharThomassen} B.~Mohar, C.~Thomassen, \emph{Graphs on Surfaces}, Johns Hopkins University Press, 2001.
\bibitem{NoschinskiAFP} L.~Noschinski, \emph{Planarity certificates}, Archive of Formal Proofs (2015).
\end{thebibliography}

\end{document}
